% Theory / Methods section -- biexciton in a strongly oblate ellipsoidal QD.
% Notation follows ../../00-theory/Notation and Coordinates.md, Single Particle.md,
% Exciton.md, Biexciton.md, and Non-Linear Optical Response.md.
% Intended to be \input into the main manuscript; requires amsmath, amssymb.

\section{Theory and Method}
\label{sec:theory}

\subsection{Model and confinement}

We consider a strongly oblate ellipsoidal quantum dot (SOEQD) grown from GaAs,
whose boundary is the semi-ellipsoid
\begin{equation}
  \mathcal{D}=\left\{(x,y,z):\;\frac{x^{2}+y^{2}}{a^{2}}+\frac{z^{2}}{c^{2}}\le 1,\;z\ge 0\right\},
  \qquad \frac{c}{a}\ll 1 ,
  \label{eq:domain}
\end{equation}
with $a$ and $c$ the large (in-plane) and small (axial) semiaxes. The carriers
are taken to be hard-wall confined,
\begin{equation}
  V_{\mathrm{conf}}(\mathbf r)=
  \begin{cases}
    0, & \mathbf r\in\mathcal{D},\\[2pt]
    \infty, & \mathbf r\notin\mathcal{D},
  \end{cases}
  \qquad \Psi\big|_{\partial\mathcal{D}}=0 .
  \label{eq:hardwall}
\end{equation}
Throughout, electrons are labelled $1,2$ and holes $a,b$; a particle position is
written in cylindrical coordinates $\mathbf r_i=(z_i,r_i,\phi_i)$, and the
interparticle separations are
\begin{equation}
  r_{ij}=\sqrt{(z_i-z_j)^{2}+r_i^{2}+r_j^{2}-2r_ir_j\cos(\phi_i-\phi_j)} .
  \label{eq:pairdist}
\end{equation}

Because the dot is strongly oblate, the axial motion is quantized on a much
larger energy scale than the in-plane motion. We therefore separate the fast
$z$ degree of freedom adiabatically at fixed in-plane radius $r$. With
$z_{\max}(r)=c\sqrt{1-r^{2}/a^{2}}$, the hard-wall axial problem yields
\begin{equation}
  \phi_\nu(z;r)=\sqrt{\frac{2}{z_{\max}}}\,
  \sin\!\left(\frac{\nu\pi z}{z_{\max}}\right),
  \qquad
  \varepsilon_\nu(r)=\nu^{2}\varepsilon_z\,\frac{1}{1-r^{2}/a^{2}},
  \qquad
  \varepsilon_z=\frac{\hbar^{2}\pi^{2}}{2m^{*}c^{2}} .
  \label{eq:zmotion}
\end{equation}
For $r\ll a$ the axial energy $\varepsilon_\nu(r)\simeq\nu^{2}\varepsilon_z(1+r^{2}/a^{2})$
acts as a parabolic in-plane potential. The carrier dynamics thus reduce to an
effective two-dimensional oscillator of frequency
$\Omega_\nu=\sqrt{2\nu^{2}\varepsilon_z/(m^{*}a^{2})}$, a result equivalent to the
quasi--two-dimensional reduction obtained by integrating out the axial
coordinate, and the in-plane Coulomb interactions inherit the planar separations
$\rho_{ij}=|\boldsymbol\rho_i-\boldsymbol\rho_j|$.

\subsection{Single-particle states}

In the adiabatic, lowest-axial-subband approximation the single-particle
spectrum and orbitals are those of a planar harmonic oscillator,
\begin{align}
  E_{\nu,n,m} &= \nu^{2}\varepsilon_z+\hbar\Omega_\nu\,(2n+|m|+1),
  \label{eq:spE}\\[4pt]
  \Psi_{\nu,n,m}(z,r,\varphi) &= \phi_\nu(z;r)\,
  \mathcal N\,r^{|m|}\,e^{-r^{2}/(2\ell_\nu^{2})}\,
  L_n^{|m|}\!\left(\frac{r^{2}}{\ell_\nu^{2}}\right)e^{im\varphi},
  \qquad \ell_\nu=\sqrt{\frac{\hbar}{m^{*}\Omega_\nu}},
  \label{eq:spPsi}
\end{align}
with radial, angular, and axial quantum numbers $n$, $m$, $\nu$, normalization
$\mathcal N$, and associated Laguerre polynomials $L_n^{|m|}$. Equations
\eqref{eq:spE}--\eqref{eq:spPsi} furnish the building blocks $\psi_e$ and
$\psi_h$ from which the correlated few-particle states are assembled. Excited
levels are generated by raising the radial quantum number, since for $c\ll a$
the axial excitations cost far more energy than the radial ones.

\subsection{Variational exciton and biexciton states}

The exciton Hamiltonian splits into a single-particle part and the
electron--hole attraction,
\begin{equation}
  H_X=h_e(1)+h_h(a)-\frac{\kappa}{r_{1a}},
  \qquad
  \kappa=\frac{e^{2}}{4\pi\varepsilon_0\varepsilon_r}.
  \label{eq:Hx}
\end{equation}
Starting from the uncorrelated product $\Psi_0=\psi_e(1)\psi_h(a)$, with
$H_{\mathrm{sp}}\Psi_0=E_0\Psi_0$ and $E_0=E_e+E_h$, electron--hole correlation
is introduced through a single Jastrow factor,
\begin{equation}
  \Psi_X=\Psi_0\,F,\qquad F=e^{-\alpha r_{1a}},
  \label{eq:Xtrial}
\end{equation}
with $\alpha$ a variational parameter. The expectation value of the energy
separates into kinetic, interaction, and norm integrals,
$E=E_0+(\mathcal K+\mathcal V)/\mathcal N$; using
$\nabla_iF=-\alpha F\,\nabla_i r_{1a}$ and $|\nabla_1 r_{1a}|^{2}=|\nabla_a r_{1a}|^{2}=1$
the kinetic term reduces to
$\mathcal K=\big(\tfrac{\hbar^{2}}{2m_e}+\tfrac{\hbar^{2}}{2m_h}\big)\alpha^{2}\mathcal N$,
so that
\begin{equation}
  E_X=E_0+\left(\frac{\hbar^{2}}{2m_e}+\frac{\hbar^{2}}{2m_h}\right)\alpha^{2}
  +\frac{\mathcal V}{\mathcal N},
  \qquad
  \mathcal V=-\kappa\!\int\!\Psi_0^{2}\,\frac{e^{-2\alpha r_{1a}}}{r_{1a}}\,d\tau .
  \label{eq:Xenergy}
\end{equation}

The biexciton is the four-particle complex of two electrons ($1,2$) and two
holes ($a,b$), with
\begin{equation}
  H_{XX}=h_e(1)+h_e(2)+h_h(a)+h_h(b)
  +\kappa\!\left(\frac{1}{r_{12}}+\frac{1}{r_{ab}}
  -\frac{1}{r_{1a}}-\frac{1}{r_{1b}}-\frac{1}{r_{2a}}-\frac{1}{r_{2b}}\right).
  \label{eq:Hxx}
\end{equation}
The trial state again multiplies the uncorrelated product
$\Psi_0=\psi_e(1)\psi_e(2)\psi_h(a)\psi_h(b)$ ($E_0=2E_e+2E_h$) by a correlation
factor that is symmetric under electron and hole exchange,
\begin{equation}
  \Psi_{XX}=\Psi_0\,F,\qquad F=G\,(P+Q),
  \label{eq:XXtrial}
\end{equation}
where the hole--hole correlation and the direct and cross electron--hole pairings are
\begin{align}
  G &= r_{ab}^{\,\gamma}\,e^{-\delta r_{ab}},\\
  P &= \exp\!\big[-\alpha(r_{1a}+r_{2b})-\beta(r_{1b}+r_{2a})\big],\\
  Q &= \exp\!\big[-\alpha(r_{1b}+r_{2a})-\beta(r_{1a}+r_{2b})\big],
  \label{eq:GPQ}
\end{align}
with four variational parameters $\alpha,\beta,\gamma,\delta$. Squaring the
correlation factor, $F^{2}=G^{2}(P^{2}+Q^{2}+2PQ)$, and using the
$a\!\leftrightarrow\! b$ (equivalently $1\!\leftrightarrow\! 2$) exchange symmetry,
the norm and interaction integrals collapse into direct ($D$) and cross ($X$)
contributions,
\begin{equation}
  \mathcal N=2\big(N^{(D)}+N^{(X)}\big),
  \qquad
  \mathcal V=2\big(V^{(D)}+V^{(X)}\big),
  \label{eq:DXdecomp}
\end{equation}
and likewise for the kinetic kernels $K_e^{(D,X)}$, $K_h^{(D,X)}$ obtained from
$\nabla_i\ln G$, $\nabla_i\ln P$, and $\nabla_i\ln Q$ together with the
pair-distance derivatives of Eq.~\eqref{eq:pairdist}. The variational biexciton
energy is then
\begin{equation}
  E_{XX}=E_0+
  \frac{\dfrac{\hbar^{2}}{m_e}\big(K_e^{(D)}+K_e^{(X)}\big)
        +\dfrac{\hbar^{2}}{m_h}\big(K_h^{(D)}+K_h^{(X)}\big)
        +V^{(D)}+V^{(X)}}
       {N^{(D)}+N^{(X)}} .
  \label{eq:XXenergy}
\end{equation}
Minimizing Eqs.~\eqref{eq:Xenergy} and \eqref{eq:XXenergy} over the variational
parameters yields the exciton and biexciton energies and wave functions as
functions of the geometrical parameters $a$ and $c$.

\subsection{Oscillator strengths and transition dipoles}

The optical response is organized around the three relevant levels: the empty-dot
ground state $g$ (no carriers), the exciton $e$, and the biexciton $b$, with
transition energies $\hbar\omega_{eg}=E_e-E_g$, $\hbar\omega_{be}=E_b-E_e$, and
$\hbar\omega_{bg}=E_b-E_g$. The exciton-creation and biexciton-creation
oscillator strengths are
\begin{equation}
  f_{eg}=\frac{2}{m_0\hbar\omega_{eg}}\,\big|\langle e|p|g\rangle\big|^{2},
  \qquad
  f_{be}=\frac{2}{m_0\hbar\omega_{be}}\,\big|\langle b|p|e\rangle\big|^{2},
  \label{eq:osc}
\end{equation}
with $m_0$ the free-electron mass and $p$ the momentum operator, and the
corresponding transition dipole moments follow from
\begin{equation}
  \mu_{ij}^{2}=\frac{\hbar e^{2}}{2m_0\,\omega_{ij}}\,f_{ij}.
  \label{eq:dipole}
\end{equation}
The radiative population-decay rate of the exciton, which sets the overall scale
of the non-linear response, is
$\Gamma_e=\dfrac{2n_r e^{2}\omega_{eg}^{2}}{3m_0 s^{3}}\,f_{eg}$, with $n_r$ the
refractive index and $s$ the speed of light.

\subsection{Third-order optical response}

Within the three-level model and in the degenerate limit of equal photon
energies $\hbar\omega_1=\hbar\omega_2\equiv\hbar\omega$, the third-order
susceptibility is~\cite{boyd}
\begin{align}
  \chi^{(3)}(\omega)
  ={}&-\frac{i|\mu_{eg}|^{4}}{2}\,
      \frac{1}{i(\hbar\omega_{eg}-2\hbar\omega)+\hbar\Gamma_{eg}}\,
      \frac{1}{\hbar\Gamma_e}
      \left(\frac{1}{i(\hbar\omega_{eg}-\hbar\omega)+\hbar\Gamma_{eg}}
      +\frac{1}{i(\hbar\omega-\hbar\omega_{eg})+\hbar\Gamma_{eg}}\right)
      \nonumber\\
  &+\frac{i|\mu_{eg}|^{2}|\mu_{be}|^{2}}{4}\,
      \frac{1}{i(\hbar\omega_{be}-2\hbar\omega)+\hbar\Gamma_{be}}\,
      \frac{1}{\hbar\Gamma_e}
      \left(\frac{1}{i(\hbar\omega_{eg}-\hbar\omega)+\hbar\Gamma_{eg}}
      +\frac{1}{i(\hbar\omega-\hbar\omega_{eg})+\hbar\Gamma_{eg}}\right)
      \nonumber\\
  &-\frac{i|\mu_{eg}|^{2}|\mu_{be}|^{2}}{4}\,
      \frac{1}{i(\hbar\omega_{eg}-2\hbar\omega)+\hbar\Gamma_{eg}}\,
      \frac{1}{i(\hbar\omega_{bg}-2\hbar\omega)+\hbar\Gamma_{bg}}\,
      \frac{1}{i(\hbar\omega_{eg}-\hbar\omega)+\hbar\Gamma_{eg}}
      \nonumber\\
  &+\frac{i|\mu_{eg}|^{2}|\mu_{be}|^{2}}{4}\,
      \frac{1}{i(\hbar\omega_{be}-2\hbar\omega)+\hbar\Gamma_{be}}\,
      \frac{1}{i(\hbar\omega_{bg}-2\hbar\omega)+\hbar\Gamma_{bg}}\,
      \frac{1}{i(\hbar\omega_{eg}-\hbar\omega)+\hbar\Gamma_{eg}} ,
  \label{eq:chi3}
\end{align}
where $\hbar\Gamma_{ij}$ is the dephasing rate of the $i\!\to\! j$ transition.
Equation~\eqref{eq:chi3} contains three spectrally distinct features. Near the
exciton one-photon resonance $\omega=\omega_{eg}$,
\begin{equation}
  \chi^{(3)}\big|_{\omega_{eg}}\simeq
  -\frac{i|\mu_{eg}|^{4}}{2\hbar\Gamma_e}\,
  \frac{2\hbar\Gamma_{eg}}{(\hbar\omega-\hbar\omega_{eg})^{2}+(\hbar\Gamma_{eg})^{2}}\,
  \frac{1}{i(\hbar\omega_{eg}-\hbar\omega)+\hbar\Gamma_{eg}},
  \label{eq:resEG}
\end{equation}
near the biexciton one-photon resonance $\omega=\omega_{be}$,
\begin{equation}
  \chi^{(3)}\big|_{\omega_{be}}\simeq
  \frac{i|\mu_{eg}|^{2}|\mu_{be}|^{2}}{4\hbar\Gamma_e}\,
  \frac{2\hbar\Gamma_{eg}}{(\hbar\omega-\hbar\omega_{eg})^{2}+(\hbar\Gamma_{eg})^{2}}\,
  \frac{1}{i(\hbar\omega_{be}-\hbar\omega)+\hbar\Gamma_{be}},
  \label{eq:resBE}
\end{equation}
and around the two-photon resonance $2\omega\approx\omega_{bg}$,
\begin{equation}
  \chi^{(3)}\big|_{2\omega_{bg}}\simeq
  \frac{2i|\mu_{eg}|^{2}|\mu_{be}|^{2}}{E_{\mathrm{bind}}^{2}}\,
  \frac{1}{i(\hbar\omega_{bg}-2\hbar\omega)+\hbar\Gamma_{bg}},
  \qquad
  E_{\mathrm{bind}}=2E_e-E_b .
  \label{eq:resBG}
\end{equation}
The real parts of $\chi^{(3)}$ change sign from negative to positive across
$\omega_{eg}$ and from positive to negative across both $\omega_{be}$ and
$\omega_{bg}$, while the imaginary part is negative at $\omega_{eg}$ and positive
at $\omega_{be}$ and $\omega_{bg}$.

\subsection{Absorption coefficients}

The total absorption coefficient is the sum of a linear and an
intensity-dependent non-linear part,
\begin{equation}
  \alpha_{\mathrm{total}}(\omega)=\alpha_0(\omega)+\alpha_2(\omega)\,I(\omega),
  \label{eq:abstot}
\end{equation}
with
\begin{equation}
  \alpha_0(\omega)=\frac{\hbar\Gamma_{eg}}
       {(\hbar\omega-\hbar\omega_{eg})^{2}+(\hbar\Gamma_{eg})^{2}},
  \qquad
  \alpha_2(\omega)=\frac{32\pi^{2}\omega}{\varepsilon_0 c^{2}}\,
       \operatorname{Im}\chi^{(3)}(\omega),
  \label{eq:abscoef}
\end{equation}
and $I(\omega)$ the incident intensity. Because the biexciton is most directly
probed through two-photon processes, we also evaluate the two-photon absorption
coefficient,
\begin{equation}
  \alpha^{(2)}(\omega)=\frac{4\pi\omega}{c\,\varepsilon_0^{1/2}}\,I(\omega)\,
  \operatorname{Im}\chi^{(3)}(\omega),
  \label{eq:twophoton}
\end{equation}
which near $2\omega\approx\omega_{bg}$ scales as
$\alpha^{(2)}\propto|\mu_{eg}|^{2}|\mu_{be}|^{2}/(E_{\mathrm{bind}}^{2}\hbar\Gamma_{bg})$,
making the two-photon channel a sensitive measure of the biexciton binding
energy and of its dependence on the dot geometry. All optical quantities are
evaluated for the GaAs material parameters $m_e^{*}=0.067\,m_0$,
$m_h^{*}=0.45\,m_0$, and $\varepsilon_r=12.91$.
